\documentclass{leveldoc}
\DocUses{foundations,slingshot,first_coil}
\DocFingerprint{9d62e5b3b275ca00}

\begin{document}

\section{What the level asks}\label{sec:asks}

The game's text:
\leveldescription

A positive particle is launched from the left edge at $(\val{xZero}, \val{yZero})$, straight
along $+x$. With nothing on the board it flies straight on and leaves through the right edge,
below the detector, the rectangle $\val{xDetLo} \le x \le \val{xDetHi}$,
$\val{yDetLo} \le y \le \val{yDetHi}$ (no acceptance window, no kinetic window, no gate). The
player places up to three magnets, each a uniformly magnetized sphere with its moment out of
the board ($\odot$) or into it ($\otimes$), of the listed strengths, inside the region
$\val{regX0} \le x \le \val{regX1}$, $\val{regY0} \le y \le \val{regY1}$. A flight ends when
the particle enters the detector, touches a magnet, or leaves the world
(\xsecref{foundations}{sec:board}).

The reference places one magnet of the largest strength, $\odot$, below the launch line near
the left edge of the region. Its field, into the board along the particle's path, bends the
particle upwards by about thirteen degrees, enough to reach the detector's left face. This
document derives the magnet's field in the plane, and finds the flight twice: by integrating the
equations of motion numerically, and by reducing them to quadratures with the two quantities
the motion conserves, the size of the momentum and the canonical angular momentum about the
magnet.

\section{Setup}\label{sec:setup}

The level file, field by field (the appendix reproduces it):
\setuptable

In symbols: the particle has charge $q = \val{q}$, mass $m = \val{m}$ and radius $0$; the speed
of light is $c = \val{c}$; the launch has kinetic energy $T_0 = \val{Tzero}$, so
$\gamma_0 = \val{gammaZero}$ and $p = \abs{\vect{p}_0} = \val{pZero}$
(\xeqref{foundations}{eq:launch}). The reference magnet is at $\vect{r}_\mathrm{m} =
(\val{xMag}, \val{yMag})$ with the strength $\mu = \val{muMag}$ (field units, $\mu_0 m/4\pi$,
\physref{2.2}); its sphere has radius $R_\mathrm{m} = \val{Rm}$. We write $q\mu = \val{qmu}$.
The world is $-1 \le x \le \val{xWorld}$, $-1 \le y \le \val{yWorld}$.

\begin{figure}[htbp]
\centering
\begin{tikzpicture}
\begin{axis}[axis equal image, width=0.92\linewidth, xmin=-1.6, xmax=31.6, ymin=-1.6, ymax=21.6,
  xtick={0,5,...,30}, ytick={0,5,...,20}, tick label style={font=\small},
  xlabel={$x$}, ylabel={$y$}, label style={font=\small}]
  \draw[black!8] (0,0) grid (30,20);
  \draw[black!45, dashed] (-1,-1) rectangle (\rawval{xWorld},\rawval{yWorld});
  \draw[black!30, dotted] (\rawval{regX0},\rawval{regY0}) rectangle (\rawval{regX1},\rawval{regY1});
  \fill[green!50!black, fill opacity=0.18] (\rawval{xDetLo},\rawval{yDetLo}) rectangle (\rawval{xDetHi},\rawval{yDetHi});
  \draw[green!40!black] (\rawval{xDetLo},\rawval{yDetLo}) rectangle (\rawval{xDetHi},\rawval{yDetHi});
  \draw[black!40, dashed] (\rawval{xMag},\rawval{yMag}) circle[radius=\rawval{rMin}];
  \draw[orange!80!black] (\rawval{xZero},\rawval{yZero}) -- (\rawval{bareX},\rawval{bareY});
  \addplot[very thick, blue!70!black] table {path.dat};
  \fill[black!70] (\rawval{xMag},\rawval{yMag}) circle[radius=\rawval{Rm}];
  \fill (\rawval{xZero},\rawval{yZero}) circle[radius=0.15];
\end{axis}
\end{tikzpicture}
\caption{The magnet (black disc), the reference flight (blue) bending up into the detector
(green), and the straight empty-board flight (orange) out through the right edge. The dashed
circle has the radius of the closest approach $r_\mathrm{min}$ \eqref{eq:rmin}: the flight
touches it and never comes nearer. The dotted rectangle is the region for magnets.}
\label{fig:board}
\end{figure}

\section{Principles and model}\label{sec:model}

\begin{principles}
  \item The Lorentz force in the relativistic equation of motion
    \xeqref{foundations}{eq:eom}; a static magnetic field does no work, so $\abs{\vect{p}}$ is
    constant (\xsecref{first_coil}{sec:nowork}).
  \item The field outside a uniformly magnetized sphere is exactly a point dipole's, of
    vector potential $\vect{A} = \vect{\mu}\times\vect{r}/r^3$ (\Jackson{§5.6, §5.10}), in the
    game's field units (\physref{2.2}).
  \item Noether's theorem: a Lagrangian invariant under rotations about an axis conserves the
    canonical angular momentum about it; for a charge, $\vect{r}\times(\vect{p} + q\vect{A})$
    (\Goldstein{Ch.~2 and 8}, \Jackson{§12.1} for the charge's Lagrangian).
  \item The board's rules (\xsecref{foundations}{sec:board}); the classical Runge--Kutta method,
    cubic Hermite interpolation and the error estimate by halving the step
    (\xsecref{foundations}{sec:numerics}), as in \xsecref{slingshot}{sec:method}.
\end{principles}

\begin{assumption}{The magnet is held in place}\label{as:fixed}
It neither moves nor turns, and its magnetization is fixed (rules of the game).
\end{assumption}

\begin{assumption}{No radiation}\label{as:radiation}
The radiated energy \xeqref{foundations}{eq:larmor} is not taken from the particle; along the
reference flight it is $\val{Erad}$, a fraction $\val{EradFrac}$ of $T_0$.
\end{assumption}

\begin{assumption}{A point particle without a moment}\label{as:point}
The particle has no magnetic moment of its own, so the field's gradient exerts no force on it
(a later level's subject).
\end{assumption}

\section{Derivation}\label{sec:derivation}

\subsection{The magnet's field in the plane}\label{sec:dipole}

The dipole's field, $\vect{B} = \nabla\times\vect{A}$, is
$\vect{B} = \mu\,(3(\uvec{m}\cdot\uvec{r})\uvec{r} - \uvec{m})/r^3$ with
$\vect{\mu} = \mu\uvec{m}$ (\Jackson{§5.6}). The moment stands perpendicular to the board,
$\uvec{m} = \uvec{z}$, so at every point of the board $\uvec{m}\cdot\uvec{r} = 0$ and
\begin{equation}\label{eq:dipole-plane}
  \vect{B} = -\frac{\mu}{r^3}\,\uvec{z}, \qquad r = \abs{\vect{x} - \vect{r}_\mathrm{m}} :
\end{equation}
perpendicular to the board, opposite to the moment, and falling as $1/r^3$. With $\mu > 0$
($\odot$) it points into the board. In the plane the vector potential is azimuthal,
$\vect{A} = (\mu/r^2)\,\uvec{\varphi}$. The force $q\vect{v}\times\vect{B}$ lies in the plane,
so the motion stays there, and
\begin{equation}\label{eq:eom-dipole}
  \deriv{\vect{x}}{t} = \vect{v} = \frac{\vect{p}}{\gamma m}, \qquad
  \deriv{\vect{p}}{t} = qB_z(\vect{x})\,(v_y, -v_x), \qquad
  B_z = -\frac{\mu}{\abs{\vect{x} - \vect{r}_\mathrm{m}}^3},
\end{equation}
from $\vect{x}_0$ and $\vect{p}_0 = (p, 0)$. At the launch the magnet is $r_0 = \val{rZero}$
away, $B_z = \val{Bzero}$, and the circle of \xeqref{first_coil}{eq:larmor-radius} would have
the radius $\val{rLzero}$: far from the magnet the path is nearly straight, and it bends where
it passes close. With $qB_z < 0$ a particle moving along $+x$ is pushed towards $+y$: the
bend is anticlockwise.

\subsection{Two conserved quantities}\label{sec:conserved}

The magnetic force does no work, so $\abs{\vect{p}} = p$ and $\gamma$ are constant
(\xsecref{first_coil}{sec:nowork}). The field \eqref{eq:dipole-plane} is unchanged by rotations
about the magnet, and so is the charge's Lagrangian $-mc^2/\gamma + q\vect{v}\cdot\vect{A}$ with
$\vect{A} = (\mu/r^2)\uvec{\varphi}$: the angle $\varphi$ about the magnet is cyclic, and its
canonical momentum is conserved,
\begin{equation}\label{eq:canonical}
  P = \big(\vect{r}\times(\vect{p} + q\vect{A})\big)_z = L + \frac{q\mu}{r},
  \qquad L = (\vect{r}\times\vect{p})_z = r\,p_\varphi ,
\end{equation}
with $\vect{r} = \vect{x} - \vect{r}_\mathrm{m}$. Directly: with
$\dot{\vect{p}} = qB_z(v_y, -v_x)$, $\dd L/\dd t = (\vect{r}\times\dot{\vect{p}})_z =
-qB_z\,(\vect{r}\cdot\vect{v}) = q\mu\,\dot r/r^2 = -\dd(q\mu/r)/\dd t$ (using
$\vect{r}\cdot\vect{v} = r\dot r$), so $L + q\mu/r$ is constant. At the launch
$L_0 = \val{Lzero}$ and $P = \val{Pcan}$.

\subsection{The flight by quadratures}\label{sec:quadrature}

With $p$ and $P$ fixed, $p_\varphi = L(r)/r = (P - q\mu/r)/r$ is known at every radius, and so
is $p_r = \pm\sqrt{p^2 - p_\varphi^2}$. The particle can only be where $\abs{L(r)} \le pr$. It
starts moving inwards ($p_r = \val{prZero}$) and turns where $\abs{L(r)} = pr$; here $L < 0$
throughout, so $P - q\mu/r = -pr$, that is
\begin{equation}\label{eq:rmin}
  p\,r^2 + P\,r - q\mu = 0, \qquad
  r_\mathrm{min} = \frac{-P + \sqrt{P^2 + 4pq\mu}}{2p} = \val{rMin} .
\end{equation}
It never comes nearer, whatever the rest of the flight. From $\dot\varphi = p_\varphi/(\gamma m
r)$ and $\dot r = p_r/(\gamma m)$, on each branch
\begin{equation}\label{eq:phi-of-r}
  \deriv{\varphi}{r} = \frac{L(r)}{r^2\sqrt{p^2 - L(r)^2/r^2}},
\end{equation}
\begin{equation}\label{eq:t-of-r}
  \deriv{t}{r} = \frac{\gamma m}{\abs{p_r}} = \frac{\gamma m}{\sqrt{p^2 - L(r)^2/r^2}} ,
\end{equation}
taken with $\dd r < 0$ on the way in and $\dd r > 0$ on the way out. So the outgoing branch at
the radius $r$ has the angle $\varphi(r) = \varphi_0 + I(r_0) + I(r)$ and the time
$t(r) = J(r_0) + J(r)$, where $I(R)$ and $J(R)$ are the integrals of \eqref{eq:phi-of-r} and
\eqref{eq:t-of-r} from $r_\mathrm{min}$ to $R$ and $\varphi_0$ is the launch's angle. Near
$r_\mathrm{min}$ both integrands grow as $1/\sqrt{r - r_\mathrm{min}}$; the substitution
$r = r_\mathrm{min} + u^2$ makes them smooth, and a standard quadrature (Gauss--Legendre or the
tanh-sinh rule, to thirty digits) evaluates them. The entry into the detector is the smallest
$r$ on the outgoing branch whose point $\vect{r}_\mathrm{m} + r(\cos\varphi, \sin\varphi)$ lies
in the detector, found by stepping $r$ outwards and bisecting; there
$\vect{p} = p_r\uvec{r} + p_\varphi\uvec{\varphi}$.

This is exact apart from the quadrature's rounding: it needs no time step. It gives the entry
at $r = \val{rEntry}$, at $t = \val{tQ}$, at $(\val{xQ}, \val{yQ})$, with momentum
$(\val{pxQ}, \val{pyQ})$.

\subsection{The method}\label{sec:method}

The march integrates \eqref{eq:eom-dipole} by the classical Runge--Kutta method
\xeqref{foundations}{eq:rk4} at the fixed step $h = \val{H}$, locates the events on the cubic
Hermite interpolant as in \xsecref{slingshot}{sec:method}, and estimates its error by halving
$h$. Along it $p$ and $P$ are the checks: they must stay constant.

\section{The reference solution by hand}\label{sec:hand}

\begin{procedure}
  \item\label{st:launch} \emph{The launch.} $\gamma_0 = \val{gammaZero}$, $p = \val{pZero}$,
    $\vect{p}_0 = (p, 0)$; $r_0 = \val{rZero}$, $L_0 = \val{Lzero}$, $P = \val{Pcan}$ by
    \eqref{eq:canonical}.
  \item\label{st:rmin} \emph{The closest approach.} \eqref{eq:rmin}:
    $r_\mathrm{min} = \val{rMin}$. The march's closest sampled point is $\val{rMinMarch}$.
  \item\label{st:march} \emph{The march.} Advance \eqref{eq:eom-dipole} by
    \xeqref{foundations}{eq:rk4}. The state every five time units (step $\val{HTable}$; the
    last column the drift of $P$):
    \begin{center}\small\val{states}\end{center}
  \item\label{st:entry} \emph{The entry.} With $h = \val{H}$, after $\val{steps}$ steps, the
    particle enters the detector through its \val{faceE} face at $t_\mathrm{e} = \val{tE}$, at
    $(\val{xE}, \val{yE})$, with momentum $(\val{pxE}, \val{pyE})$, $\val{angE}$ degrees above
    $+x$. The step-halving estimate of its error is $\val{err}$.
  \item\label{st:checks} \emph{The checks.} Along the march $\abs{\vect{p}}$ drifts by at most
    $\val{pDev}$ and $P$ by $\val{PDev}$. The quadratures of \S\ref{sec:quadrature} give the
    same end state to $\val{quadDiff}$.
  \item\label{st:rad} \emph{The radiation.} Simpson's rule on the steps gives $\val{Erad}$;
    against the step $2h$ its error is a relative $\val{EradErr}$.
\end{procedure}

\begin{result}[The reference flight]
The particle enters the detector through its \val{faceE} face at $(\val{xE}, \val{yE})$ at
$t_\mathrm{e} = \val{tE}$, with momentum $(\val{pxE}, \val{pyE})$. The level is solved.
\end{result}

\section{Comparison with the game}\label{sec:comparison}

\comparisontable

The reference flight agrees to $\val{cmpYDiff}$ in $y_\mathrm{e}$ and $\val{cmpTDiff}$ in
$t_\mathrm{e}$, the empty board to $\val{cmpBareTDiff}$: the size of the game's verified
accuracy (about $10^{-10}$, Foundations §6), above our step estimate $\val{err}$ and the
quadratures' agreement $\val{quadDiff}$. The tolerance is $10^{-9}$ (ten times the larger of
the two computations' error estimates), or the game's own preview--verification gap where that
is larger ($\val{gapRef}$ on the reference: not larger), fixed before comparing.

\section{Other solutions}\label{sec:others}

The empty board, the second placement in \code{placements.json}, does not solve the level: the
particle flies straight and leaves through the \val{bareFace} edge at
$(\val{bareX}, \val{bareY})$ at $t = \val{bareT}$. A magnet bends the particle only near it:
the bend's strength goes as $\mu$ over the closest approach cubed, roughly, so it must sit
near the path. One $\odot$ magnet below the path (or $\otimes$ above it) bends it upwards; the
reference uses the strongest, below the path at the region's edge. Several weaker magnets
along the path can share the bend. Equation \eqref{eq:rmin} gives each single-magnet
placement's closest approach without any integration, a quick test of whether the particle
would strike the magnet; a map of all solutions would take the game's verdicts, which this
document does not reproduce.

\section{What the game computes}\label{sec:engine}

The game takes a magnet's field outside its sphere as the exact point-dipole field
(\physref{2.2}), which is \eqref{eq:dipole-plane} in the plane, with exactly zero in-plane
components, and integrates \eqref{eq:eom-dipole} with its adaptive eighth-order Runge--Kutta
method (\physref{5}), event location (\physref{6}) and verification (\physref{7}). The game
does not use the conserved $P$; it holds to it to its integration accuracy, as the march
here does to $\val{PDev}$. Radiation is integrated, not applied (\physref{8}).

\section{Further reading}\label{sec:reading}

The dipole's field and a magnetized sphere: \Jackson{§5.6, §5.10}. The canonical momentum and
cyclic coordinates: \Goldstein{Ch.~2 and 8}; the charge's Lagrangian: \Jackson{§12.1}. The
forbidden regions of a charge in a dipole's field, of which \eqref{eq:rmin} is the case of a
plane, are Störmer's problem. Cyclotron motion and the magnetic force's work:
\levelref{first_coil}.

\levelfileappendix

\end{document}
